\section{与前几期的连接：Agent、Context 与语言}

这期与前面几期形成一条很清晰的应用层主线。EP139 讲 Agent 技术史，强调工具调用、边界消融和社会辐射；EP135 讲 Elys 和赛博分身，强调 context 的获取、流动与人格化网络；EP133 则从世界模型角度提醒语言之外还有物理世界。EP130 的位置，是把这些技术概念落回产品经理的操作台。

\subsection{同一个词，不同层面的含义}

\begin{longtable}{p{0.22\textwidth}p{0.34\textwidth}p{0.35\textwidth}}
\toprule
概念 & 在前几期中的侧重 & 在 EP130 中的产品含义 \\
\midrule
Agent & 工具调用、任务执行、系统边界变化 & 从 PPT/创作表达切入，形成可交互生产力对象。 \\
Context & 个人信息、社交网络、赛博分身的数据流 & 产品经理要设计模型可见上下文，而不是只画流程。 \\
语言 & EP133 中被讨论为不够覆盖物理世界 & EP130 中被视为 Agent 的治理内核和用户意图入口。 \\
AI 朋友 & EP135 更偏关系和社交网络 & EP130 强调先提供实用价值，再逐渐形成朋友感。 \\
\bottomrule
\end{longtable}

\begin{knowledgebox}{为什么这条连接重要}
如果只看 EP130，它像产品方法论；放进整个队列，它是在回答“Agent 怎么变成普通人会长期使用的产品”。技术史、context 流动、语言边界和产品心智，是同一个问题的不同层。
\end{knowledgebox}

\subsection{本章小结}

EP130 把 Agent 从技术叙事带回产品叙事：不只问模型能不能做，还问用户为什么记住它、为什么回不去、为什么愿意让它进入自己的工作和关系网络。

